\section{Bài toán và cách chia việc}

Hệ thống nhận một câu hỏi bằng tiếng Anh, tìm trong \NumCorpusChunks{} đoạn văn
(cắt từ \NumCorpusArticles{} bài báo CNN) những đoạn nhiều khả năng chứa câu trả
lời nhất, rồi để một LLM viết đáp án \emph{chỉ dựa trên} các đoạn đó, kèm số
trích dẫn \texttt{[n]} trỏ về đoạn đã dùng.

Dự án tách làm hai giai đoạn vì hai nửa của đường ống hỏng theo hai cách khác
nhau và phải đo bằng hai loại thước khác nhau.

\begin{center}\small
\begin{tabularx}{0.94\textwidth}{@{}p{2.8cm} L L@{}}
\toprule
 & \textbf{Phase 1 --- lấy đoạn văn} & \textbf{Phase 2 --- hỏi LLM} \\
\midrule
Hỏng kiểu gì & Không tìm ra đoạn chứa đáp án & Tìm ra rồi nhưng trả lời sai, hoặc bịa \\
Đo bằng gì & Hit@k, nDCG@5, MRR@5 & Answer Correctness, Faithfulness, Citation \\
Chi phí & GPU, không gọi API & Tốn tiền API mỗi lần chạy \\
Tái lập & Có, nhờ chỉ mục cố định & Chỉ trong giới hạn nhiễu của LLM \\
\bottomrule
\end{tabularx}
\end{center}

Thứ tự là bắt buộc: Phase~1 phải khóa xong trước khi Phase~2 bắt đầu. Nếu vừa
đổi cách lấy đoạn vừa đổi cách hỏi thì không có cách nào quy được cải thiện về
nguyên nhân nào. Cụ thể hơn, Phase~2 chạy lại trên một \textbf{vết truy xuất đã
đóng băng}: cùng một danh sách đoạn đã xếp hạng cho mọi cấu hình sinh, nên mọi
khác biệt quan sát được đều thuộc về tầng sinh.

\subsection{Dữ liệu và những gì EDA phát hiện}

Giai đoạn khảo sát dữ liệu không phải bước trang trí --- nó viết sẵn giả thuyết
cho từng vòng thí nghiệm về sau, và đặt ra ngưỡng phán quyết.

\begin{center}\small
\begin{tabularx}{0.94\textwidth}{@{}p{5.4cm} L@{}}
\toprule
\textbf{Phát hiện} & \textbf{Hệ quả lên thiết kế thực nghiệm} \\
\midrule
Nhãn nhiễu \NumNoiseFloor{}--\NumNoiseCeiling{} &
  Chênh lệch dưới \NumNoiseFloor{} không được coi là thắng \\
Trung vị \NumDistractorMedian{} đoạn đối thủ cùng chủ đề cho mỗi câu hỏi &
  Reranker là bắt buộc, không phải tùy chọn \\
Câu hỏi neo vào tên riêng, ngày tháng, con số &
  Dự đoán sparse thắng dense, \emph{trước} khi đo \\
\NumCorpusRestored{} bài bị cắt mất phần đuôi &
  Dựng lại kho: \NumCorpusChunksVone{} $\rightarrow$ \NumCorpusChunks{} đoạn \\
Notebook cũ đặt \texttt{JUDGE\_MODEL = GENERATOR\_MODEL} &
  Tách judge sang nhà cung cấp khác generator \\
\bottomrule
\end{tabularx}
\end{center}

\paragraph{Biên độ nhiễu \NumNoiseFloor{}.}
EDA ước lượng được rằng \NumNoiseFloor{}--\NumNoiseCeiling{} số câu trong tập
đánh giá mang nhãn có vấn đề --- đáp án chuẩn không đầy đủ, hoặc có nhiều span
đúng nhưng chỉ một được ghi nhận. Hệ quả trực tiếp: một cấu hình hơn cấu hình
khác dưới \NumNoiseFloor{} thì khoảng cách đó \emph{có thể} hoàn toàn do nhiễu
nhãn. Toàn bộ báo cáo áp dụng luật phán quyết hai điều kiện:

\keypoint{Một cấu hình chỉ được tuyên là thắng khi \textbf{(1)} khoảng cách vượt
biên độ nhiễu \NumNoiseFloor{}, \emph{và} \textbf{(2)} khoảng tin cậy 95\% của
hiệu \emph{theo cặp} không chứa 0. Có ý nghĩa thống kê là điều kiện cần, không
phải điều kiện đủ.}

\paragraph{Hai cách đọc cùng một tập câu hỏi.}
Nhiều câu hỏi NewsQA không đứng độc lập được --- chúng được viết như đang nói
tiếp một cuộc hội thoại (\emph{``How many were killed?''}). Tỉ lệ câu phụ thuộc
ngữ cảnh giảm từ 57{,}5\% xuống 11{,}3\% sau khi giải quyết tham chiếu. Bản
\texttt{original} \textbf{thiên vị sparse}: câu hỏi càng cụt thì khớp từ vựng
càng có lợi thế giả tạo. Vì hệ thống thật luôn nhận câu hỏi đầy đủ ngữ cảnh,
mọi số chính trong báo cáo đo trên tập \texttt{resolved}; \texttt{original} chỉ
dùng làm cận dưới.

\paragraph{Bootstrap theo cặp, gom theo bài báo.}
Nhiều câu hỏi cùng trỏ về một bài báo, nên chúng không phải quan sát độc lập.
Mọi khoảng tin cậy trong báo cáo tính bằng bootstrap phần trăm \emph{theo cụm
bài báo}: mỗi lần lấy mẫu bốc lại nguyên một bài kèm toàn bộ câu hỏi của nó
(\NumBootstrapSamples{} mẫu, seed 42). Coi từng câu là một lần bốc độc lập sẽ
làm khoảng tin cậy hẹp hơn thực tế.
